\iftestbuild % TEST-ONLY-BEGIN
\setcounter{section}{2}\setcounter{theorem}{34}
\fi % TEST-ONLY-END
% Source: book pp. 154--162 (PDF 168--176); no solutions (manual stops at 5B).
% \axtag{ax:N.M}: an equation that carries an Axler number (shares the item
% counter).  \exnote{...}: the small italic comment that follows some exercises.
\DeclareRobustCommand\axtag[1]{\refstepcounter{theorem}\tag{\thetheorem}\label{#1}}
\section{Upper-Triangular Matrices}

In Chapter~3 we defined the matrix of a linear map from a finite-dimensional
vector space to another finite-dimensional vector space. That matrix depends on a
choice of basis of each of the two vector spaces. Now that we are studying
operators, which map a vector space to itself, the emphasis is on using only one
basis.

\begin{definition}[Matrix of an operator, $\Mat(T)$]\label{ax:5.35}
Suppose $T\in\Lin(V)$. The \emph{matrix} of $T$ with respect to a basis
$v_1,\dots,v_n$ of $V$ is the $n$-by-$n$ matrix
\[ \Mat(T) = \begin{pmatrix} A_{1,1} & \cdots & A_{1,n}\\ \vdots & & \vdots\\
   A_{n,1} & \cdots & A_{n,n} \end{pmatrix} \]
whose entries $A_{j,k}$ are defined by
\[ Tv_k = A_{1,k}v_1+\cdots+A_{n,k}v_n. \]
The notation $\Mat\bigl(T,(v_1,\dots,v_n)\bigr)$ is used if the basis is not
clear from the context.
\end{definition}

Operators have square matrices (meaning that the number of rows equals the
number of columns), rather than the more general rectangular matrices that we
considered earlier for linear maps.

\marginnote{The $k^{\text{th}}$ column of the matrix $\Mat(T)$ is formed from the
coefficients used to write $Tv_k$ as a linear combination of the basis
$v_1,\dots,v_n$.}%
If $T$ is an operator on $\F^n$ and no basis is specified, assume that the basis
in question is the standard one (where the $k^{\text{th}}$ basis vector is $1$ in
the $k^{\text{th}}$ slot and $0$ in all other slots). You can then think of the
$k^{\text{th}}$ column of $\Mat(T)$ as $T$ applied to the $k^{\text{th}}$ basis
vector, where we identify $n$-by-$1$ column vectors with elements of $\F^n$.

\begin{example}[Matrix of an operator with respect to standard basis]\label{ax:5.36}
Define $T\in\Lin(\F^3)$ by $T(x,y,z)=(2x+y,5y+3z,8z)$. Then the matrix of $T$
with respect to the standard basis of $\F^3$ is
\[ \Mat(T) = \begin{pmatrix} 2&1&0\\ 0&5&3\\ 0&0&8 \end{pmatrix}, \]
as you should verify.
\end{example}

A central goal of linear algebra is to show that given an operator $T$ on a
finite-dimensional vector space $V$, there exists a basis of $V$ with respect to
which $T$ has a reasonably simple matrix. To make this vague formulation a bit
more precise, we might try to choose a basis of $V$ such that $\Mat(T)$ has many
$0$'s.

If $V$ is a finite-dimensional complex vector space, then we already know enough
to show that there is a basis of $V$ with respect to which the matrix of $T$ has
$0$'s everywhere in the first column, except possibly the first entry. In other
words, there is a basis of $V$ with respect to which the matrix of $T$ looks like
\[ \begin{pmatrix} \lambda & & & \\ 0 & & * & \\ \vdots & & & \\ 0 & & &
   \end{pmatrix}; \]
here $*$ denotes the entries in all columns other than the first column. To prove
this, let $\lambda$ be an eigenvalue of $T$ (one exists by \ref{ax:5.19}) and let
$v$ be a corresponding eigenvector. Extend $v$ to a basis of $V$. Then the matrix
of $T$ with respect to this basis has the form above. Soon we will see that we
can choose a basis of $V$ with respect to which the matrix of $T$ has even more
$0$'s.

\begin{definition}[Diagonal of a matrix]\label{ax:5.37}
The \emph{diagonal} of a square matrix consists of the entries on the line from
the upper left corner to the bottom right corner.
\end{definition}

For example, the diagonal of the matrix
\[ \Mat(T) = \begin{pmatrix} \textcolor{red}{2}&1&0\\ 0&\textcolor{red}{5}&3\\
   0&0&\textcolor{red}{8} \end{pmatrix} \]
from Example~\ref{ax:5.36} consists of the entries $2,5,8$, which are shown in
red in the matrix above.

\begin{definition}[Upper-triangular matrix]\label{ax:5.38}
A square matrix is called \emph{upper triangular} if all entries below the
diagonal are $0$.
\end{definition}

For example, the $3$-by-$3$ matrix above is upper triangular.

Typically we represent an upper-triangular matrix in the form
\[ \begin{pmatrix} \lambda_1 & & *\\ & \ddots & \\ 0 & & \lambda_n
   \end{pmatrix}; \]
\marginnote{We often use $*$ to denote matrix entries that we do not know or
that are irrelevant to the questions being discussed.}%
the $0$ in the matrix above indicates that all entries below the diagonal in this
$n$-by-$n$ matrix equal $0$. Upper-triangular matrices can be considered
reasonably simple---if $n$ is large, then at least almost half the entries in an
$n$-by-$n$ upper-triangular matrix are $0$.

The next result provides a useful connection between upper-triangular matrices
and invariant subspaces.

\begin{result}[Conditions for upper-triangular matrix]\label{ax:5.39}
Suppose $T\in\Lin(V)$ and $v_1,\dots,v_n$ is a basis of $V$. Then the following
are equivalent.
\begin{parts}
\item The matrix of $T$ with respect to $v_1,\dots,v_n$ is upper triangular.
\item $\spn(v_1,\dots,v_k)$ is invariant under $T$ for each $k=1,\dots,n$.
\item $Tv_k\in\spn(v_1,\dots,v_k)$ for each $k=1,\dots,n$.
\end{parts}
\end{result}

\begin{proof}
First suppose (a) holds. To prove that (b) holds, suppose $k\in\{1,\dots,n\}$.
If $j\in\{1,\dots,n\}$, then
\[ Tv_j\in\spn(v_1,\dots,v_j) \]
because the matrix of $T$ with respect to $v_1,\dots,v_n$ is upper triangular.
Because $\spn(v_1,\dots,v_j)\subseteq\spn(v_1,\dots,v_k)$ if $j\le k$, we see
that
\[ Tv_j\in\spn(v_1,\dots,v_k) \]
for each $j\in\{1,\dots,k\}$. Thus $\spn(v_1,\dots,v_k)$ is invariant under $T$,
completing the proof that (a) implies (b).

Now suppose (b) holds, so $\spn(v_1,\dots,v_k)$ is invariant under $T$ for each
$k=1,\dots,n$. In particular, $Tv_k\in\spn(v_1,\dots,v_k)$ for each
$k=1,\dots,n$. Thus (b) implies (c).

Now suppose (c) holds, so $Tv_k\in\spn(v_1,\dots,v_k)$ for each $k=1,\dots,n$.
This means that when writing each $Tv_k$ as a linear combination of the basis
vectors $v_1,\dots,v_n$, we need to use only the vectors $v_1,\dots,v_k$. Hence
all entries under the diagonal of $\Mat(T)$ are $0$. Thus $\Mat(T)$ is an
upper-triangular matrix, completing the proof that (c) implies (a).

We have shown that (a) $\implies$ (b) $\implies$ (c) $\implies$ (a), which shows
that (a), (b), and (c) are equivalent.
\end{proof}

The next result tells us that if $T\in\Lin(V)$ and with respect to some basis of
$V$ we have
\[ \Mat(T) = \begin{pmatrix} \lambda_1 & & *\\ & \ddots & \\ 0 & & \lambda_n
   \end{pmatrix}, \]
then $T$ satisfies a simple equation depending on $\lambda_1,\dots,\lambda_n$.

\begin{result}[Equation satisfied by operator with upper-triangular matrix]%
\label{ax:5.40}
Suppose $T\in\Lin(V)$ and $V$ has a basis with respect to which $T$ has an
upper-triangular matrix with diagonal entries $\lambda_1,\dots,\lambda_n$. Then
\[ (T-\lambda_1I)\cdots(T-\lambda_nI)=0. \]
\end{result}

\begin{proof}
Let $v_1,\dots,v_n$ denote a basis of $V$ with respect to which $T$ has an
upper-triangular matrix with diagonal entries $\lambda_1,\dots,\lambda_n$. Then
$Tv_1=\lambda_1v_1$, which means that $(T-\lambda_1I)v_1=0$, which implies that
$(T-\lambda_1I)\cdots(T-\lambda_mI)v_1=0$ for $m=1,\dots,n$ (using the
commutativity of each $T-\lambda_jI$ with each $T-\lambda_kI$).

Note that $(T-\lambda_2I)v_2\in\spn(v_1)$. Thus
$(T-\lambda_1I)(T-\lambda_2I)v_2=0$ (by the previous paragraph), which implies
that $(T-\lambda_1I)\cdots(T-\lambda_mI)v_2=0$ for $m=2,\dots,n$ (using the
commutativity of each $T-\lambda_jI$ with each $T-\lambda_kI$).

Note that $(T-\lambda_3I)v_3\in\spn(v_1,v_2)$. Thus by the previous paragraph,
$(T-\lambda_1I)(T-\lambda_2I)(T-\lambda_3I)v_3=0$, which implies that
$(T-\lambda_1I)\cdots(T-\lambda_mI)v_3=0$ for $m=3,\dots,n$ (using the
commutativity of each $T-\lambda_jI$ with each $T-\lambda_kI$).

Continuing this pattern, we see that $(T-\lambda_1I)\cdots(T-\lambda_nI)v_k=0$
for each $k=1,\dots,n$. Thus $(T-\lambda_1I)\cdots(T-\lambda_nI)$ is the $0$
operator because it is $0$ on each vector in a basis of $V$.
\end{proof}

Unfortunately no method exists for exactly computing the eigenvalues of an
operator from its matrix. However, if we are fortunate enough to find a basis with
respect to which the matrix of the operator is upper triangular, then the problem
of computing the eigenvalues becomes trivial, as the next result shows.

\begin{result}[Determination of eigenvalues from upper-triangular matrix]%
\label{ax:5.41}
Suppose $T\in\Lin(V)$ has an upper-triangular matrix with respect to some basis
of $V$. Then the eigenvalues of $T$ are precisely the entries on the diagonal of
that upper-triangular matrix.
\end{result}

\begin{proof}
Suppose $v_1,\dots,v_n$ is a basis of $V$ with respect to which $T$ has an
upper-triangular matrix
\[ \Mat(T) = \begin{pmatrix} \lambda_1 & & *\\ & \ddots & \\ 0 & & \lambda_n
   \end{pmatrix}. \]

Because $Tv_1=\lambda_1v_1$, we see that $\lambda_1$ is an eigenvalue of $T$.

Suppose $k\in\{2,\dots,n\}$. Then $(T-\lambda_kI)v_k\in\spn(v_1,\dots,v_{k-1})$.
Thus $T-\lambda_kI$ maps $\spn(v_1,\dots,v_k)$ into $\spn(v_1,\dots,v_{k-1})$.
Because
\[ \dim\spn(v_1,\dots,v_k)=k \quad\text{and}\quad
   \dim\spn(v_1,\dots,v_{k-1})=k-1, \]
this implies that $T-\lambda_kI$ restricted to $\spn(v_1,\dots,v_k)$ is not
injective (by \ref{ax:3.22}). Thus there exists $v\in\spn(v_1,\dots,v_k)$ such
that $v\ne0$ and $(T-\lambda_kI)v=0$. Thus $\lambda_k$ is an eigenvalue of $T$.
Hence we have shown that every entry on the diagonal of $\Mat(T)$ is an
eigenvalue of $T$.

To prove $T$ has no other eigenvalues, let $q$ be the polynomial defined by
$q(z)=(z-\lambda_1)\cdots(z-\lambda_n)$. Then $q(T)=0$ (by \ref{ax:5.40}). Hence
$q$ is a polynomial multiple of the minimal polynomial of $T$ (by
\ref{ax:5.29}). Thus every zero of the minimal polynomial of $T$ is a zero of
$q$. Because the zeros of the minimal polynomial of $T$ are the eigenvalues of
$T$ (by \ref{ax:5.27}), this implies that every eigenvalue of $T$ is a zero of
$q$. Hence the eigenvalues of $T$ are all contained in the list
$\lambda_1,\dots,\lambda_n$.
\end{proof}

\begin{example}[Eigenvalues via an upper-triangular matrix]\label{ax:5.42}
Define $T\in\Lin(\F^3)$ by $T(x,y,z)=(2x+y,5y+3z,8z)$. The matrix of $T$ with
respect to the standard basis is
\[ \Mat(T) = \begin{pmatrix} 2&1&0\\ 0&5&3\\ 0&0&8 \end{pmatrix}. \]
Now \ref{ax:5.41} implies that the eigenvalues of $T$ are $2$, $5$, and $8$.
\end{example}

The next example illustrates \ref{ax:5.44}: an operator has an upper-triangular
matrix with respect to some basis if and only if the minimal polynomial of the
operator is the product of polynomials of degree $1$.

\begin{example}[Whether $T$ has an upper-triangular matrix can depend on $\F$]%
\label{ax:5.43}
Define $T\in\Lin(\F^4)$ by
\[ T(z_1,z_2,z_3,z_4)=(-z_2,z_1,2z_1+3z_3,z_3+3z_4). \]
Thus with respect to the standard basis of $\F^4$, the matrix of $T$ is
\[ \begin{pmatrix} 0&-1&0&0\\ 1&0&0&0\\ 2&0&3&0\\ 0&0&1&3 \end{pmatrix}. \]
You can ask a computer to verify that the minimal polynomial of $T$ is the
polynomial $p$ defined by
\[ p(z)=9-6z+10z^2-6z^3+z^4. \]

First consider the case $\F=\R$. Then the polynomial $p$ factors as
\[ p(z)=(z^2+1)(z-3)(z-3), \]
with no further factorization of $z^2+1$ as the product of two polynomials of
degree $1$ with real coefficients. Thus \ref{ax:5.44} states that there does not
exist a basis of $\R^4$ with respect to which $T$ has an upper-triangular matrix.

Now consider the case $\F=\C$. Then the polynomial $p$ factors as
\[ p(z)=(z-i)(z+i)(z-3)(z-3), \]
where all factors above have the form $z-\lambda_k$. Thus \ref{ax:5.44} states
that there is a basis of $\C^4$ with respect to which $T$ has an
upper-triangular matrix. Indeed, you can verify that with respect to the basis
$(4-3i,-3-4i,-3+i,1)$, $(4+3i,-3+4i,-3-i,1)$, $(0,0,0,1)$, $(0,0,1,0)$ of $\C^4$,
the operator $T$ has the upper-triangular matrix
\[ \begin{pmatrix} i&0&0&0\\ 0&-i&0&0\\ 0&0&3&1\\ 0&0&0&3 \end{pmatrix}. \]
\end{example}

\begin{result}[Necessary and sufficient condition to have an upper-triangular
matrix]\label{ax:5.44}
Suppose $V$ is finite-dimensional and $T\in\Lin(V)$. Then $T$ has an
upper-triangular matrix with respect to some basis of $V$ if and only if the
minimal polynomial of $T$ equals $(z-\lambda_1)\cdots(z-\lambda_m)$ for some
$\lambda_1,\dots,\lambda_m\in\F$.
\end{result}

\begin{proof}
First suppose $T$ has an upper-triangular matrix with respect to some basis of
$V$. Let $\alpha_1,\dots,\alpha_n$ denote the diagonal entries of that matrix.
Define a polynomial $q\in\Poly(\F)$ by
\[ q(z)=(z-\alpha_1)\cdots(z-\alpha_n). \]
Then $q(T)=0$, by \ref{ax:5.40}. Hence $q$ is a polynomial multiple of the
minimal polynomial of $T$, by \ref{ax:5.29}. Thus the minimal polynomial of $T$
equals $(z-\lambda_1)\cdots(z-\lambda_m)$ for some
$\lambda_1,\dots,\lambda_m\in\F$ with
$\{\lambda_1,\dots,\lambda_m\}\subseteq\{\alpha_1,\dots,\alpha_n\}$.

To prove the implication in the other direction, now suppose the minimal
polynomial of $T$ equals $(z-\lambda_1)\cdots(z-\lambda_m)$ for some
$\lambda_1,\dots,\lambda_m\in\F$. We will use induction on $m$. To get started,
if $m=1$ then $z-\lambda_1$ is the minimal polynomial of $T$, which implies that
$T=\lambda_1I$, which implies that the matrix of $T$ (with respect to any basis
of $V$) is upper triangular.

Now suppose $m>1$ and the desired result holds for all smaller positive integers.
Let
\[ U=\range(T-\lambda_mI). \]
Then $U$ is invariant under $T$ [this is a special case of \ref{ax:5.18} with
$p(z)=z-\lambda_m$]. Thus $T|_U$ is an operator on $U$.

If $u\in U$, then $u=(T-\lambda_mI)v$ for some $v\in V$ and
\[ (T-\lambda_1I)\cdots(T-\lambda_{m-1}I)u=(T-\lambda_1I)\cdots(T-\lambda_mI)v=0. \]
Hence $(z-\lambda_1)\cdots(z-\lambda_{m-1})$ is a polynomial multiple of the
minimal polynomial of $T|_U$, by \ref{ax:5.29}. Thus the minimal polynomial of
$T|_U$ is the product of at most $m-1$ terms of the form $z-\lambda_k$.

By our induction hypothesis, there is a basis $u_1,\dots,u_M$ of $U$ with respect
to which $T|_U$ has an upper-triangular matrix. Thus for each
$k\in\{1,\dots,M\}$, we have (using \ref{ax:5.39})
\[ Tu_k=(T|_U)(u_k)\in\spn(u_1,\dots,u_k). \axtag{ax:5.45} \]

Extend $u_1,\dots,u_M$ to a basis $u_1,\dots,u_M,v_1,\dots,v_N$ of $V$. If
$k\in\{1,\dots,N\}$, then
\[ Tv_k=(T-\lambda_mI)v_k+\lambda_mv_k. \]
The definition of $U$ shows that $(T-\lambda_mI)v_k\in U=\spn(u_1,\dots,u_M)$.
Thus the equation above shows that
\[ Tv_k\in\spn(u_1,\dots,u_M,v_1,\dots,v_k). \axtag{ax:5.46} \]

From \ref{ax:5.45} and \ref{ax:5.46}, we conclude (using \ref{ax:5.39}) that $T$
has an upper-triangular matrix with respect to the basis
$u_1,\dots,u_M,v_1,\dots,v_N$ of $V$, as desired.
\end{proof}

The set of numbers $\{\lambda_1,\dots,\lambda_m\}$ from the previous result
equals the set of eigenvalues of $T$ (because the set of zeros of the minimal
polynomial of $T$ equals the set of eigenvalues of $T$, by \ref{ax:5.27}),
although the list $\lambda_1,\dots,\lambda_m$ in the previous result may contain
repetitions.

In Chapter~8 we will improve even the wonderful result below; see \ref{ax:8.37}
and \ref{ax:8.46}.

\begin{result}[If $\F=\C$, then every operator on $V$ has an upper-triangular
matrix]\label{ax:5.47}
Suppose $V$ is a finite-dimensional complex vector space and $T\in\Lin(V)$. Then
$T$ has an upper-triangular matrix with respect to some basis of $V$.
\end{result}

\begin{proof}
The desired result follows immediately from \ref{ax:5.44} and the second version
of the fundamental theorem of algebra (see \ref{ax:4.13}).
\end{proof}

For an extension of the result above to two operators $S$ and $T$ such that
\[ ST=TS, \]
see \ref{ax:5.80}. Also, for an extension to more than two operators, see
Exercise~9(b) in Section~5E.

\textbf{Caution:} If an operator $T\in\Lin(V)$ has an upper-triangular matrix
with respect to some basis $v_1,\dots,v_n$ of $V$, then the eigenvalues of $T$
are exactly the entries on the diagonal of $\Mat(T)$, as shown by
\ref{ax:5.41}, and furthermore $v_1$ is an eigenvector of $T$. However,
$v_2,\dots,v_n$ need not be eigenvectors of $T$. Indeed, a basis vector $v_k$ is
an eigenvector of $T$ if and only if all entries in the $k^{\text{th}}$ column
of the matrix of $T$ are $0$, except possibly the $k^{\text{th}}$ entry.

\marginnote{The row echelon form of the matrix of an operator does not give us a
list of the eigenvalues of the operator. In contrast, an upper-triangular matrix
with respect to some basis gives us a list of all the eigenvalues of the
operator. However, there is no method for computing exactly such an
upper-triangular matrix, even though \ref{ax:5.47} guarantees its existence if
$\F=\C$.}%
You may recall from a previous course that every matrix of numbers can be
changed to a matrix in what is called row echelon form. If one begins with a
square matrix, the matrix in row echelon form will be an upper-triangular matrix.
Do not confuse this upper-triangular matrix with the upper-triangular matrix of
an operator with respect to some basis whose existence is proclaimed by
\ref{ax:5.47} (if $\F=\C$)---there is no connection between these
upper-triangular matrices.

% ------------------------------------------------------------------ exercises
\begin{exc}

\begin{exercise}
Prove or give a counterexample: If $T\in\Lin(V)$ and $T^2$ has an
upper-triangular matrix with respect to some basis of $V$, then $T$ has an
upper-triangular matrix with respect to some basis of $V$.
\end{exercise}

\begin{exercise}
Suppose $A$ and $B$ are upper-triangular matrices of the same size, with
$\alpha_1,\dots,\alpha_n$ on the diagonal of $A$ and $\beta_1,\dots,\beta_n$ on
the diagonal of $B$.
\begin{parts}
\item Show that $A+B$ is an upper-triangular matrix with
  $\alpha_1+\beta_1,\dots,\alpha_n+\beta_n$ on the diagonal.
\item Show that $AB$ is an upper-triangular matrix with
  $\alpha_1\beta_1,\dots,\alpha_n\beta_n$ on the diagonal.
\end{parts}
\exnote{The results in this exercise are used in the proof of \ref{ax:5.81}.}
\end{exercise}

\begin{exercise}
Suppose $T\in\Lin(V)$ is invertible and $v_1,\dots,v_n$ is a basis of $V$ with
respect to which the matrix of $T$ is upper triangular, with
$\lambda_1,\dots,\lambda_n$ on the diagonal. Show that the matrix of $T^{-1}$ is
also upper triangular with respect to the basis $v_1,\dots,v_n$, with
\[ \frac{1}{\lambda_1},\dots,\frac{1}{\lambda_n} \]
on the diagonal.
\end{exercise}

\begin{exercise}
Give an example of an operator whose matrix with respect to some basis contains
only $0$'s on the diagonal, but the operator is invertible.
\exnote{This exercise and the exercise below show that \ref{ax:5.41} fails
without the hypothesis that an upper-triangular matrix is under consideration.}
\end{exercise}

\begin{exercise}
Give an example of an operator whose matrix with respect to some basis contains
only nonzero numbers on the diagonal, but the operator is not invertible.
\end{exercise}

\begin{exercise}
Suppose $\F=\C$, $V$ is finite-dimensional, and $T\in\Lin(V)$. Prove that if
$k\in\{1,\dots,\dim V\}$, then $V$ has a $k$-dimensional subspace invariant under
$T$.
\end{exercise}

\begin{exercise}
Suppose $V$ is finite-dimensional, $T\in\Lin(V)$, and $v\in V$.
\begin{parts}
\item Prove that there exists a unique monic polynomial $p_v$ of smallest degree
  such that $p_v(T)v=0$.
\item Prove that the minimal polynomial of $T$ is a polynomial multiple of $p_v$.
\end{parts}
\end{exercise}

\begin{exercise}
Suppose $V$ is finite-dimensional, $T\in\Lin(V)$, and there exists a nonzero
vector $v\in V$ such that $T^2v+2Tv=-2v$.
\begin{parts}
\item Prove that if $\F=\R$, then there does not exist a basis of $V$ with
  respect to which $T$ has an upper-triangular matrix.
\item Prove that if $\F=\C$ and $A$ is an upper-triangular matrix that equals the
  matrix of $T$ with respect to some basis of $V$, then $-1+i$ or $-1-i$ appears
  on the diagonal of $A$.
\end{parts}
\end{exercise}

\begin{exercise}
Suppose $B$ is a square matrix with complex entries. Prove that there exists an
invertible square matrix $A$ with complex entries such that $A^{-1}BA$ is an
upper-triangular matrix.
\end{exercise}

\begin{exercise}
Suppose $T\in\Lin(V)$ and $v_1,\dots,v_n$ is a basis of $V$. Show that the
following are equivalent.
\begin{parts}
\item The matrix of $T$ with respect to $v_1,\dots,v_n$ is lower triangular.
\item $\spn(v_k,\dots,v_n)$ is invariant under $T$ for each $k=1,\dots,n$.
\item $Tv_k\in\spn(v_k,\dots,v_n)$ for each $k=1,\dots,n$.
\end{parts}
\exnote{A square matrix is called \textbf{lower triangular} if all entries above
the diagonal are $0$.}
\end{exercise}

\begin{exercise}
Suppose $\F=\C$ and $V$ is finite-dimensional. Prove that if $T\in\Lin(V)$, then
there exists a basis of $V$ with respect to which $T$ has a lower-triangular
matrix.
\end{exercise}

\begin{exercise}
Suppose $V$ is finite-dimensional, $T\in\Lin(V)$ has an upper-triangular matrix
with respect to some basis of $V$, and $U$ is a subspace of $V$ that is invariant
under $T$.
\begin{parts}
\item Prove that $T|_U$ has an upper-triangular matrix with respect to some basis
  of $U$.
\item Prove that the quotient operator $T/U$ has an upper-triangular matrix with
  respect to some basis of $V/U$.
\end{parts}
\exnote{The quotient operator $T/U$ was defined in Exercise~38 in Section~5A.}
\end{exercise}

\begin{exercise}
Suppose $V$ is finite-dimensional and $T\in\Lin(V)$. Suppose there exists a
subspace $U$ of $V$ that is invariant under $T$ such that $T|_U$ has an
upper-triangular matrix with respect to some basis of $U$ and also $T/U$ has an
upper-triangular matrix with respect to some basis of $V/U$. Prove that $T$ has
an upper-triangular matrix with respect to some basis of $V$.
\end{exercise}

\begin{exercise}
Suppose $V$ is finite-dimensional and $T\in\Lin(V)$. Prove that $T$ has an
upper-triangular matrix with respect to some basis of $V$ if and only if the dual
operator $T'$ has an upper-triangular matrix with respect to some basis of the
dual space $V'$.
\end{exercise}
\end{exc}
